%!TEX root = ../Thesis.tex

% % Neue Befehle
\newcommand{\HRule}[2]{\noindent\rule[#1]{\linewidth}{#2}} % Horiz. Linie
\newcommand{\vlinespace}[1]{\vspace*{#1\baselineskip}} % Abstand
\newcommand{\titleemph}[1]{\textbf{#1}} % Hervorheben

\begin{titlepage}
 \sffamily % Titelseite in seriefenloser Schrift
      \includegraphics[scale=0.7]{ressources/img/logo_itdesigners}
      \hfill
      % Logo Hochschule
      \includegraphics[width=5cm]{ressources/img/logo_hfu}
      \HRule{13pt}{1pt}
      \centering
      \Large
      \vlinespace{1}\\
      \workTyp\\
      \vspace{1em}
      \huge
      \workTitel\\
%
      \Large
      \vlinespace{2}
          im Master-Studiengang \workStudiengang\\
          der Hochschule Furtwangen\\

      \vlinespace{2}
      \workNameStudent
%
   \vfill
   \raggedright
%
   \large
   \titleemph{Zeitraum:} \workSemester \\ % Nur bei Abschluss-Arbeiten
   \titleemph{Prüfer:} \workPruefer \\
   \titleemph{Zweitprüfer:} \workBetreuer \\ % Nur bei Abschluss-Arbeiten

 % Folgenden Abschnitt nur bei Industrie-Arbeiten darstellen
   \vlinespace{1}
   \HRule{13pt}{1pt} \\
   \titleemph{Firma:} \workFirma \\
   \titleemph{Betreuer:} \workBetreuer
%
\end{titlepage}
